\chapter{Toehold-mediated strand displacement}
An invader binds a short exposed segment of a DNA complex. Does the incumbent
strand necessarily leave? No. Binding may reverse before the branch point has
crossed the duplex. The useful computation is not the arrival of an invader;
it is a change in which strand occupies a substrate, followed by an observable
output. This chapter reconstructs that distinction from molecular geometry,
through first-passage probability, to finite material and readout budgets.

Prerequisites are polarity and complementarity from Chapter 5, conservation
from Chapter 6, stochastic rates from Chapter 7, and the separation of logical
and physical state from Chapter 17. We use domain-level molecular drawings and
assigned-rate models. No sequence-specific design or new experiment is claimed.

\section{Specify the strands before writing a reaction}
Name an invader $I$, an incumbent output $O$, and a substrate $S$. A domain is
a contiguous segment treated as one recognition unit in a coarse description.
Let $t$ denote an exposed \Term{toehold} and $b$ a branch-migration domain.
Write the invader from $5'$ to $3'$ as $tb$. The aligned substrate is written
from $3'$ to $5'$ as $t^*b^*$, where the star denotes complementary pairing
in this alignment. If that substrate is instead listed in its own $5'$ to
$3'$ sequence direction, its domain order reverses. Do not order oligonucleotides
by copying an aligned lower strand as if it had the upper strand's polarity.

Initially $O$ occupies $b^*$, leaving $t^*$ exposed. The invader can bind that
exposed region and compete for the adjacent paired region. Backbone continuity
does not change in this elementary exchange. The process substitutes pairing
partners; it does not require cutting the substrate or synthesizing a new strand.
Srinivas and colleagues distinguish reversible toehold binding, branch migration,
and incumbent release in their mechanistic treatment \cite{srinivas}.

\Plate{01-exchange}{One substrate, three pairing configurations. Solid colored
strokes are covalent backbones; short vertical rungs indicate paired regions.
The invader and incumbent have the same $5'$ to $3'$ orientation and compete
for an antiparallel substrate. Released output retains its backbone. Geometry
is schematic, with domain lengths not to scale.}{fig:exchange}

The \Term{branch point} marks the boundary between substrate bases paired to
the invader and those still paired to the incumbent. Its movement can reverse.
A domain-level arrow suppresses many such exchanges; it is not a claim that
the whole domain flips in one elementary event. Neither correct domain names
nor Watson--Crick complementarity rules alone determine a kinetic rate.

\section{Make the logical reaction conserve physical strands}
Let $G=S{:}O$ be the initial gate complex and $W=S{:}I$ the spent complex.
An ideal net reaction is
\begin{equation}
 I+G\longrightarrow W+O.
 \label{eq:net}
\end{equation}
The two complexes each contain two strands. Initially the system has one
invader strand and a two-strand gate; finally it has a two-strand spent complex
and one output strand. Three strands remain three. Two reactant \emph{species}
do not imply two physical strands.

In a closed, well-mixed ideal batch with reaction extent $x$ in concentration
units and no initial products,
\begin{equation}
 [I]=I_0-x,\quad [G]=G_0-x,\quad [W]=[O]=x.
 \label{eq:inventory}
\end{equation}
Hence $[I]+[W]=I_0$, $[G]+[W]=G_0$, and $[G]+[O]=G_0$.
These track invader, substrate and incumbent strands, respectively. They are
stronger diagnostics than merely checking that concentrations stay nonnegative.
A product model can conserve a sum while silently assigning the wrong strand
identity to that sum.

\Plate{04-inventory}{The three strand inventories cross the reaction boundary
unchanged. Gate consumption limits output even if input is abundant. The output
can trigger a subsequent gate only if it has that gate's required recognition
interface; a labeled logical wire is not automatically a compatible strand.}{fig:inventory}

This gate is not catalytic in $I$: the invader remains in $W$. A catalytic
design needs additional reactions that release or regenerate the input.
Likewise, feeding one released strand into two downstream gates is a competition
for material, not free Boolean fan-out. Additional fuel, copy production, or
a different signal encoding must be accounted for explicitly.

\section{Turn a bound complex into a first-passage problem}
Let $m\geq1$ count idealized branch-migration steps, and let $j$ be the number
already completed. The transient states are $j=0,\ldots,m-1$; $P$ is released
product and $D$ is detached invader. Starting at $j=0$ means the toehold is
already bound. We make four modeling choices: forward rate $k_+>0$, backward
rate $k_-\geq0$ for $j>0$, detachment rate $d\geq0$ only at $j=0$, and an
absorbing product when the forward step leaves $m-1$. Rates have units
$\mathrm{s}^{-1}$.

These choices are a teaching reduction. Real interfaces can have additional
dissociation paths, sequence-dependent barriers, and reversible product binding.
In particular, treating product as absorbing is an observation-window or
reaction-design assumption, not a universal thermodynamic statement.

\Plate{02-walk}{The state graph distinguishes detachment from productive
release. Backward motion does not necessarily detach the invader: it may return
to the toehold-bound state. The first-passage calculation ends at either absorbing
outcome; repeated binding attempts are a separate process.}{fig:walk}

Let $Q$ be the transient row generator. Its off-diagonal entry $Q_{ij}$ is
the rate from $i$ to $j$; its diagonal is minus the sum of \emph{all} outgoing
rates, including absorption. Let $r_P$ collect direct rates to $P$. The
detachment exit is $d$ in row zero and zero elsewhere. Thus
$Q\boldsymbol 1+r_P+r_D=0$. A generator with row sum zero on the transient
subspace would have forgotten absorption.
\Listing{branch_chain}

\section{Derive success probability, not just a mean rate}
Let $p_j$ be the probability of release before detachment from state $j$.
Condition on the first event. At an interior point with total exit rate $a_j$,
\[
 p_j=\sum_{i\ne j}\frac{Q_{ji}}{a_j}p_i+\frac{(r_P)_j}{a_j}.
\]
Multiplying by $a_j$ and using $Q_{jj}=-a_j$ gives the backward equation
\begin{equation}
 -Qp=r_P.\label{eq:probability}
\end{equation}
This is a linear boundary-value problem, not a Monte Carlo estimate.

\begin{proposition}[Uniform migration with detachment at the origin]
If $k_+=k_-=k>0$, then
\begin{equation}
 p_j=\frac{1+jd/k}{1+md/k},\qquad p_0=\frac{1}{1+md/k}.
 \label{eq:closed}
\end{equation}
\end{proposition}
\begin{proof}
Interior first-step equations give $p_{j+1}-2p_j+p_{j-1}=0$, so successive
differences are constant and $p_j=A+Bj$. At zero,
$k(p_1-p_0)-dp_0=0$, hence $B=dA/k$. The released boundary is $p_m=1$,
which gives $A(1+md/k)=1$ and the stated expression.
\end{proof}
When $d=0$, success is certain in this finite model. When $d/k$ grows,
most bound attempts detach. Larger $m$ lowers success at fixed rates, but this
does not predict a universal physical length law: the rates and available
paths may also change when a real sequence changes.

\Plate{03-commitment}{Computed release-before-detachment probabilities for
two migration lengths, at $k=2\ \mathrm{s}^{-1}$. The graph is a parameter
sensitivity analysis of Equation~\ref{eq:closed}, not a calibration against
toehold length or an experimental yield curve.}{fig:commitment}

\section{Distinguish all attempts from successful attempts}
Let $\tau_j$ be expected time to either absorbing outcome. First-step analysis
includes the holding time $1/a_j$, giving
\begin{equation}
 -Q\tau=\boldsymbol 1.
\end{equation}
The mean time to \emph{successful} release, conditional on success, is different.
Define $u_j=\mathbb E_j[T\,1_{\{P\}}]$. Since the eventual success probability
from the state occupied at time $t$ is $p$, integration over surviving paths
gives $u=\int_0^\infty e^{Qt}p\,dt$, hence
\begin{equation}
 -Qu=p,\qquad \mathbb E_j[T\mid P]=u_j/p_j.
\end{equation}
All transient states here reach absorption almost surely and $p_j>0$ because
the forward path has positive rates. Other models may have traps or zero
success probabilities; conditional means then require separate handling.
\Listing{first_passage}
\begin{center}\input{\Generated/passage.tex}\end{center}
The generated table uses $k_+=k_-=2$ and $d=1$, all in inverse seconds.
A fast average failed attempt must not be reported as a fast successful
computation. As an independent limit check, setting $d=0$ yields
$\tau_0=m(m+1)/(2k)$ for the reflecting-origin, absorbing-end uniform walk.
The quadratic growth follows by summing the successive differences of the
mean-time recurrence, even though a direct forward path has only $m$ steps.

The dense reference solver uses $O(m^3)$ arithmetic and $O(m^2)$ storage;
the tridiagonal structure permits linear-time elimination. Dense linear algebra
is retained for inspection and testing, not offered as a scalable molecular
simulator. Chapter 28 treats reusable solvers and numerical diagnostics.

\section{Connect microscopic attempts to a coarse reaction rate}
Let $k_{\rm on}$ have units $\mathrm{M}^{-1}\mathrm{s}^{-1}$ and let free
invader concentration be $c$ M. Binding to one free gate then has hazard
$a=k_{\rm on}c$. If each failed encounter restores the same free gate and
the bath remains constant, the renewal reasoning of Chapter 17 gives
\begin{equation}
 \mathbb E[T_{\rm product}]=\frac{1/a+\tau_0}{p_0}.
 \label{eq:renewal}
\end{equation}
The waiting time to bind and the bound residence time both count. When
$a\tau_0\ll1$, bound encounters occupy little of the total time and an
effective second-order description uses $k_{\rm eff}\approx k_{\rm on}p_0$.
Without that separation, replacing the pathway by one instantaneous reaction
can give the right eventual branch fraction and the wrong latency.

The physical literature identifies an initiation penalty and a distinction
between branch migration and simple end fraying \cite{srinivas}. Our constant
rates omit those details deliberately. For sequence-resolved prediction, one
needs an appropriate energy/rate model, calibrated conditions and validation
data. Stronger binding cannot be assumed to accelerate every later step.
Toehold exchange can also make the net process reversible; Zhang and Winfree's
work treats that architecture explicitly \cite{zhang}. It is not identical to
the absorbing-end construction solved here.

\section{Solve a finite-pool batch without creating material}
Assume Equation~\ref{eq:net} is irreversible on the observation interval,
well mixed, and accurately represented by constant $k_{\rm eff}$. Then
\begin{equation}
 \frac{dx}{dt}=k_{\rm eff}(I_0-x)(G_0-x),\qquad x(0)=0.
\end{equation}
If $I_0=G_0=c_0$, separation of variables gives
$1/(c_0-x)-1/c_0=k_{\rm eff}t$ and therefore
\begin{equation}
 x(t)=c_0\frac{k_{\rm eff}c_0t}{1+k_{\rm eff}c_0t}.
 \label{eq:batch}
\end{equation}
For unequal pools, let $H=\max(I_0,G_0)$, $L=\min(I_0,G_0)$ and
$\Delta=H-L>0$. Integrating partial fractions gives
\begin{equation}
 x(t)=\frac{LH(1-e^{-k_{\rm eff}\Delta t})}
 {\Delta+L(1-e^{-k_{\rm eff}\Delta t})}.
\end{equation}
Both formulas approach $L$, never more than the limiting reagent. The reference
uses \texttt{expm1} to resolve a small exponential difference and handles equal
pools separately. Here concentrations are in nM, so the numerical rate is in
$\mathrm{nM}^{-1}\mathrm{s}^{-1}$; multiply that rate by $10^9$ to express it
in $\mathrm{M}^{-1}\mathrm{s}^{-1}$.
\Listing{conversion}
\begin{center}\input{\Generated/batch.tex}\end{center}
The table assigns $I_0=G_0=100$ nM and $k_{\rm eff}=10^{-4}$
$\mathrm{nM}^{-1}\mathrm{s}^{-1}$. At 100 seconds only half the gate is
converted. Treating input as permanently fixed at its initial concentration
would predict faster conversion and hide depletion. This is why the following
excess-input readout model is labeled as a different approximation.

\section{Leakage creates a decision window, not only an error bar}
Define an input-present conversion hazard $\lambda$ and a background leak
hazard $\ell$, both inverse seconds. Assume excess input keeps $\lambda$
constant, reporters perfectly reflect converted gate, and every gate converts
at most once. Then normalized signals satisfy
\begin{equation}
 y_1(t)=1-e^{-(\lambda+\ell)t},\qquad
 y_0(t)=1-e^{-\ell t}.
\end{equation}
Both tend to one for $\ell>0$. Waiting indefinitely destroys separation in
this model. For threshold $0<\theta<1$, a noiseless discrimination window is
\begin{equation}
 \frac{-\log(1-\theta)}{\lambda+\ell}\leq t
 <\frac{-\log(1-\theta)}{\ell}.
\end{equation}
With $\lambda=0.01$, $\ell=0.001$ and $\theta=0.5$, the window is approximately
63.0 to 693.1 seconds. These are model outputs, not recommended assay timings.
For zero leak the upper limit becomes infinite; for zero input hazard no
strict window exists. Noise, incomplete reporters and uncertain rates shrink
or eliminate the practically usable window.
\Plate{05-leak}{Assigned-hazard signals cross the same threshold at different
times. The dashed blank curve eventually produces a false positive under this
decision rule. A valid signal comparison needs a declared readout time and
matched background controls, not only an attractive final endpoint.}{fig:leak}
\Listing{reporter}

Leakage is a mechanism class, not one universal constant. Unintended exposure,
alternative pairing, impurities and reporter reactions can produce different
time courses. A professional experiment distinguishes no-input background,
noncomplementary input, a disabled intended interaction and a known output
standard. The last control tests measurement independently of computation.
An intensity increase alone does not identify which reaction produced it.

\section{Composition, research frontier and evidence}
Local correctness does not guarantee that a cascade is correct. A released
output must match the next gate, remain available rather than sequestered,
and arrive before competing reactions dominate. Preparing multiple species
changes the interaction set; testing each gate alone does not test crosstalk
in the mixture. A useful interface contract specifies strand identities,
orientation, signal encoding, allowable concentration range, consumed resources,
and how output is distinguished from exposed domains on intermediate complexes.

The modern research question is not merely whether a strand can displace
another. It is whether a compiled molecular system preserves intended dynamics
under simultaneous reactions and changing inventories. Soloveichik, Seelig and
Winfree's reaction-network construction gives a theoretical bridge between
formal chemical dynamics and strand-displacement implementations \cite{crn}.
Its existence does not validate arbitrary sequences or remove its operating
assumptions. Chapter 20 develops stoichiometric and compositional semantics
before Chapter 21 addresses circuit-level restoration and fan-out.

A September 2026 contrast is the Scaffolded DNA Computer of Stérin and
colleagues \cite{sterin}. Its design favors a target equilibrium configuration
and allows replacement pathways rather than enforcing one detailed kinetic
route. That is a different computational contract from our absorbing-end walk.
The comparison raises a useful research question: which requirements truly
need kinetic-path control, and which can be enforced by a global energy
landscape? Neither architecture licenses assuming that every favorable endpoint
is reached quickly. We do not transfer its reported experimental performance
to the toy pathway or to arbitrary strand-displacement circuits.

For an AI-assisted design workflow, retain candidate provenance and explicit
constraints, use a physical model to reject invalid candidates, rank uncertainty
as well as predicted performance, and reserve independent measurements for
validation. Model-generated sequences are hypotheses. A high predicted score
cannot replace material accounting, specificity tests or experimental evidence.
This principle connects molecular programming to Evolutor without equating
software gating with a biochemical reaction.

\section{Exercises with worked answers}
\begin{enumerate}
\item Which strand must expose the complementary toehold? \textbf{Answer:}
The substrate in $G$, here $t^*$ aligned $3'$ to $5'$. The invader carries $t$.
\item Is branch migration backbone cutting? \textbf{Answer:} No; base-pair
partners change while the three covalent strands remain continuous.
\item With $m=4$, $k=2$ and $d=1$, compute success. \textbf{Answer:}
$p_0=1/(1+4/2)=1/3$ per bound attempt, not per starting free molecule.
\item Double both $k$ and $d$. \textbf{Answer:} Success is unchanged because
$d/k$ is unchanged; all first-passage times halve.
\item Derive the boundary condition at zero. \textbf{Answer:}
$p_0=[k/(k+d)]p_1+[d/(k+d)]0$, so $k(p_1-p_0)=dp_0$.
\item Why is $-Q\tau=1$ different from $-Qp=r_P$? \textbf{Answer:}
The first rewards time in transient states; the second rewards only the
product exit. Their right-hand sides have different units and meanings.
\item At $d=0$, $m=4$, $k=2$, find the mean. \textbf{Answer:}
$4\cdot5/(2\cdot2)=5$ seconds; it is not simply $m/k=2$ seconds.
\item Can 30 nM input and 100 nM gate produce 100 nM output here?
\textbf{Answer:} No. Without input regeneration, $x\leq30$ nM.
\item What are the units of $k_{\rm eff}c_0t$? \textbf{Answer:} Dimensionless,
provided the rate and concentration use the same molarity scale.
\item At equal pools, solve for the half-conversion time. \textbf{Answer:}
Equation~\ref{eq:batch} gives $t_{1/2}=1/(k_{\rm eff}c_0)$.
\item Why can later measurement be worse? \textbf{Answer:} Under positive
leak, blank conversion approaches full output and eventually crosses threshold.
\item Design a falsification test for a supposed leak-free cascade.
\textbf{Rubric:} State the hypothesis and detection limit; include no-input,
noncomplementary input, individual-gate and full-mixture conditions; independently
calibrate reporter response; predeclare time windows and replicates; quantify
uncertainty. No observed signal supports an upper bound, not proof of zero leak.
\end{enumerate}

\section{Vocabulary and next question}
\Term{First passage}: the first arrival at a specified outcome set.
\Term{Transient generator}: rate matrix restricted to states before absorption.
\Term{Reaction extent}: how much net reaction has occurred in declared units.
\Term{Leakage}: output-producing behavior outside the intended input pathway.
\Term{Interface contract}: the molecular and measurement assumptions required
for a component to compose with another. The next chapter turns such contracts
into chemical reaction networks while retaining their material meaning.
